
\section{Related Work}
\label{sec:related}

PV forecasting studies provide the technical background for the proposed workflow. Reviews of photovoltaic power forecasting emphasize that useful models depend on forecast horizon, input availability, sky variability, and the evaluation protocol \cite{AntonanzasPVForecastingReview2016,YangSolarForecastVerification2020}. Clear-sky-index and clearness-index transformations are widely used to separate deterministic solar geometry from atmospheric attenuation, especially when irradiance or PV output has strong daily and seasonal structure \cite{LauretBenchmarkingML2015,LauretClearSkyIndex2022}. Solar-resource variability studies further show why pointwise error, ramp behavior, and temporal aggregation should be inspected together rather than judged by one headline metric \cite{PerezVariability2016}.

Gradient-boosted tree models are attractive for tabular PV forecasting because they handle nonlinear feature interactions without requiring a large neural sequence model. XGBoost is a standard scalable implementation of regularized gradient tree boosting \cite{ChenXGBoost2016}. Recent hybrid solar-forecasting work also combines tree-based learning with physics-informed features and temporal correction models, which supports the broader direction of using physical variables to structure data-driven predictors without claiming that the base learning algorithm is new \cite{PhysicsInformedXGBLSTM2026}. Solar-position and clear-sky modeling tools such as pvlib provide the reference context for more detailed solar-geometry implementations, even though the present experiment uses a simplified, reproducible proxy \cite{HolmgrenPvlib2018}.

Microgrid and building energy-management literature provides the main application context for this work. Practical microgrid operation requires forecasting because PV output affects storage dispatch, load scheduling, and grid-interaction decisions \cite{ref01}. Building-sector energy-management reviews similarly emphasize that sensing, data quality, and control strategy must be treated together rather than as isolated modeling tasks \cite{ref02}. Comparative reviews of microgrid energy-management strategies and optimization techniques show that renewable uncertainty remains a central difficulty for operational scheduling \cite{ref05,ref06,ref11}.

Data-driven methods are now common in smart buildings and distributed-energy systems because they can learn nonlinear relationships between weather, user behavior, and local generation. Deep-learning and data-driven energy-management studies show how measured building and resource data can support predictive control and scheduling under uncertainty \cite{ref04,ref07,ref08,ref09}. In residential and community settings, home energy-management systems, multi-home coordination, and community microgrid designs show why short-interval PV forecasts matter for operational decision making \cite{ref10,ref18,ref20,ref25}.

Demand-response research is also relevant because forecast errors can shift when flexible loads or storage should be activated. Coordinated rail-transit response, PV intermittency management, multi-generation demand response, and low-voltage distribution control all depend on knowing when renewable generation will be available or uncertain \cite{ref03,ref13,ref16,ref21}. Price-based and data-driven demand-response studies extend this point to residential areas, where uncertainty in user behavior and local PV production affects scheduling quality \cite{ref22,ref23,ref24}.

Optimization and intelligent-control studies show a broad set of tools for handling renewable variability, including isolated multi-microgrid management, hybrid demand-response optimization, incentive-based demand response, smart-market management, transactive energy, reinforcement learning, and decentralized demand-side management \cite{ref14,ref15,ref28,ref29,ref31,ref32,ref38}. These papers motivate accurate PV forecasting as an enabling input, but the present manuscript remains focused on plant-level PV power prediction rather than on solving the downstream optimization problem.

The proposed method differs from the energy-management literature in scope. It does not introduce a new market mechanism, demand-response algorithm, or microgrid controller. Instead, it builds an operational PV forecasting workflow whose outputs can later be used by the scheduling and energy-management systems surveyed in the ordered references \cite{ref35,ref36,ref37,ref39,ref40}.
